% SOURCE_AUTHORITY: sga5_fr_workpass.tex, lines 8665--8693 (LF-normalized unit).
% SOURCE_SHA256: 791F4EFFC5E02832D5D77ED03518C8156D6F07E4C8238B03545DB93D883FBB28
\subsubsection*{1.4.6}
Seja \(F\) um \(\mathbb Z_\ell\)-feixe construtível localmente livre (1.3.3). O funtor composto de \(p\) e o funtor \(G\longrightarrow\cHom(F,G)\) (1.3.3) torna inversíveis as multiplicações pelos elementos de \(\mathbb Z_\ell\) e define, portanto, um funtor
\[
\cHom(F,\cdot):\mathbb Q_\ell\text{-}\fc(X)\longrightarrow \mathbb Q_\ell\text{-}\fc(X).
\]
Se \(F\) e \(F'\) são dois \(\mathbb Z_\ell\)-feixes localmente livres construtíveis que definem \(\mathbb Q_\ell\)-feixes isomorfos, existe um morfismo de \(\mathbb Z_\ell\)-feixes
\[
u:F\longrightarrow F'
\]
cujo núcleo e conúcleo são de torção. Isso mostra imediatamente que os funtores \(\cHom(F,\cdot)\) e \(\cHom(F',\cdot)\) de \(\mathbb Q_\ell\text{-}\fc(X)\) em si mesma são isomorfos.

Seja agora \(F\) um \(\mathbb Q_\ell\)-feixe constante torcido construtível. Se se denota também por \(F\) o \(\mathbb Z_\ell\)-feixe correspondente, da interpretação das restrições de \(F\) às componentes conexas \(Y\) de \(X\) como \(\pi_1(Y)\)-módulos contínuos segue que \(F\) se escreve, de maneira única salvo isomorfismo, na forma
\[
F=P\oplus T,
\]
com \(P\) localmente livre e \(T\) de torção. Isso permite definir, para todo \(\mathbb Q_\ell\)-feixe \(G\), um \(\mathbb Q_\ell\)-feixe
\[
\cHom(F,G)=\cHom(P,G).
\]
Um isomorfismo \(u:F\to F'\) de \(\mathbb Q_\ell\)-feixes define, para todo \(G\), um isomorfismo
\[
\cHom(F',G)\xrightarrow{\sim}\cHom(F,G).
\]
Mais precisamente, define-se assim um bifuntor
\[
\cHom:(\mathbb Q_\ell\text{-}\fct(X))^\circ\times \mathbb Q_\ell\text{-}\fc(X)
\longrightarrow \mathbb Q_\ell\text{-}\fc(X),
\]
onde \(\mathbb Q_\ell\text{-}\fct(X)\) denota a categoria dos \(\mathbb Q_\ell\)-feixes constantes torcidos construtíveis. Este bifuntor possui as propriedades de exatidão habituais.
